% =====================================================================
%  ch6_theta.tex — §6 theta 函数
% =====================================================================
\section{Theta 函数}\label{ss:theta}

下面说明构造的各个环节如何组装出镜像函子 $\cF(M)\to \Perf(X_K)$ 的具体描述。所考虑的特殊情形是 Floer 复形 $CF^*((L,\cE),(L',\cE'))$ 集中在次数零，此时给出显式映射
\begin{equation}\label{eq:theta}
\Phi_{L,L'}:CF^0((L,\cE),(L',\cE'))\to \Hom(\Phi(L,\cE),\Phi(L',\cE')).
\end{equation}

考虑两个对象 $(L,\cE)$、$(L',\cE')$，以及次数零的交点 $x\in L\cap L'$，设 $x\in P_v$。下面说明如何给它配置映射
\begin{equation}\label{eq:map_for_intersection_points}
\Phi_{L,L',x}:\Hom_{\cO_K}(\cE_x,\cE'_x)\to \Hom_{\cO_{U_v}}(\Phi(L,\cE)_{\mid U_v},\Phi(L',\cE')_{\mid U_v}).
\end{equation}

取图 $G$ 中包含 $x\in M_e$ 的半边 $e/v$。记 $r_{e/v}(x)\in \Z$ 为 $\phi_H^1(L)$ 沿从 $p_v$ 到 $x$ 的路径关于 $L'$ 的负方向旋转数；记 $S_{e/v}(x)$ 为 $M_e$ 内连接 $x$ 与 $p_v$ 的圆盘带符号面积，等价地，是 $L$ 与 $L'$ 在 $(0,\tau(x))\times S^1\subset M_e$ 的万有覆盖上围出的区域之面积（取当 $\tau\to\tau(x)$ 时彼此靠近的提升）。当 $x=p_v$ 时有 $r_{e/v}(p_v)=0$、$S_{e/v}(p_v)=0$。若 $e$ 连接 $v$ 与 $v'$，则
\begin{equation}\label{eq:relations_on_r_and_S}
r_{e/v}(x)+r_{e/v'}(x)=r_{e}(L,L'),\quad
S_{e/v}(x)-S_{e/v'}(x)+r_{e/v'}(x)A_e=S_{e/v}(L,L').
\end{equation}

对每个这样的 $x$ 配置单项式 $T^{S_{e/v}(x)}t_{e/v}^{-r_{e/v}(x)}$，看作 $U_v$ 上的函数。再按如下公式定义映射~\eqref{eq:map_for_intersection_points}：
\begin{equation}\label{eq:monomial_for_a_point}
\Phi_{L,L',x}(\varphi)=(R_{\cE',x,v}\varphi R_{\cE,v,x})\otimes T^{S_{e/v}(x)}t_{e/v}^{-r_{e/v}(x)}\in \Hom_{\cO_{U_v}}(\Phi(L,\cE)_{\mid U_v},\Phi(L',\cE')_{\mid U_v}),
\end{equation}
其中 $\varphi\in\Hom(\cE_x,\cE'_x)$，$R_{\cE,v,x}$、$R_{\cE',x,v}$ 记单值性。此外，两个态射
\[
(R_{\cE',x,v}\varphi R_{\cE,v,x})\otimes T^{S_{e/v}(x)}t_{e/v}^{-r_{e/v}(x)},\qquad
(R_{\cE',x,v'}\varphi R_{\cE,v',x})\otimes T^{S_{e/v'}(x)}t_{e/v'}^{-r_{e/v'}(x)}
\]
在 $U_e$ 上一致。这可由 $\Phi(L,\cE)$、$\Phi(L',\cE')$ 的黏合数据~\eqref{eq:gluingdata}，以及利用式~\eqref{eq:relations_on_r_and_S} 的观察
\begin{multline*}
\bigl(T^{-S_{e/v}(L,L')}t_{e/v}^{r_e(L,L')}\bigr)\bigl(T^{S_{e/v}(x)}t_{e/v}^{-r_{e/v}(x)}\bigr)
=T^{S_{e/v}(x)-S_{e/v}(L,L')}t_{e/v}^{r_{e/v'}(x)}\\
=T^{S_{e/v}(x)-S_{e/v}(L,L')+r_{e/v'}(x)A_e}t_{e/v'}^{-r_{e/v'}(x)}
=T^{S_{e/v'}(x)}t_{e/v'}^{-r_{e/v'}(x)}
\end{multline*}
推出。

现在引入一些记号。取 $G$ 中任一约化路径，写作 $\gamma=(v_0,e_1,\dots,e_n,v_n)$。它给出映射 $g_{\gamma}:Y_{v_0}\to Y_{v_n}$，
\[
g_{\gamma}=g_{e_n/v_{n-1}}\circ\dots\circ g_{e_1/v_0}.
\]
记黏合态射为
\[
u_{L,L',\gamma}:\Hom(\cE_{v_0},\cE'_{v_0}\otimes \cO_{Y_{v_0}-F_{v_0}})\to g_{\gamma}^*\bigl(\Hom(\cE_{v_n},\cE'_{v_n})\otimes \cO_{Y_{v_n}-F_{v_n}}\bigr).
\]

态射~\eqref{eq:theta} 由对态射~\eqref{eq:monomial_for_a_point}「取平均」给出。即对半边 $e_0/v_0$、点 $x\in L\cap L'\cap (\mathrm{int}(M_{e_0})\sqcup\{v_0\})$ 与态射 $\varphi:\cE_x\to\cE'_x$，对任意顶点 $v\in V(G)$ 令
\begin{equation}\label{eq:averaging}
\Phi_{L,L'}(\varphi)_{\mid U_v}=\sum_{\gamma:v_0\to v}g_{\gamma*}\Bigl(u_{L,L',\gamma}\bigl((R_{\cE',x,v_0}\varphi R_{\cE,v_0,x})\otimes T^{S_{e_0/v_0}(x)}t_{e_0/v_0}^{-r_{e_0/v_0}(x)}\bigr)\Bigr).
\end{equation}
此和收敛，因为已假定 Floer 复形 $CF^{*}((L,\cE),(L',\cE'))$ 集中在次数零。$\Phi_{L,L'}(\varphi)$ 到不同 $U_v$ 的限制在交上一致，故得到良定义的向量丛态射 $\Phi(L,\cE)\to \Phi(L',\cE')$。

\begin{example}\label{ex:elliptic}
取图 $G$ 只有一个顶点 $v$、一条连接 $v$ 与自身的内部边 $e$、以及一条外部边 $e'$。于是 $M$ 是 $M_e=\PP^1/(0\sim\infty)$ 与 $M_{e'}=\C$ 在 $M_e$ 的节点处粘成之并。对顶点 $v$ 选取组合数据，使与边 $e$ 两端相应的坐标 $t^{\pm 1}$ 互为倒数，并在第三个穿孔点取值为 $1$。

此时镜像是带穿孔的椭圆曲线 $X_K=(K^*-q_e^\Z)/q_e^\Z$，其中 $q_e=T^{A_e}$。取 v.b. 型对象 $L_0,L_1$，使得 $L_1$ 从 $M_e$ 的节点出发绕行一圈回到自身时相对 $L_0$ 顺时针旋转一次，即 $r_e(L_0,L_1)=1$，且 $S_{e/v}(L_0,L_1)=\frac12 A_e$；两者都配以平凡局部系统。按第~\ref{ss:assignbundles} 节的构造，镜像函子把 $L_0$ 送到结构层 $\cO_X$，把 $L_1$ 送到 $(K^*-q_e^\Z)$ 上平凡线丛的某个商。

对绕边 $e$ 走 $n$ 圈、从 $v$ 回到自身的路径 $\gamma_n$，坐标变换 $g_{\gamma_n}$ 由 $t\mapsto q_e^{-n} t$ 给出；线丛 $\Phi(L_1)$ 的纤维（等价地，$\Hom(\Phi(L_0),\Phi(L_1))$）相应的黏合映射 $u_{L_1,\gamma_n}$（即 $u_{L_0,L_1,\gamma_n}$）由乘以 $q_e^{-n^2/2} t^{n}$ 给出。

取 $x=p_v$ 为位于节点处的 $L_0$ 与 $L_1$ 之交，$\varphi\in CF^0(L_0,L_1)$ 为此交点与平凡局部系统间恒等映射所定义的态射。则式~\eqref{eq:averaging} 变为
\[
\Phi_{L_0,L_1}(\varphi)=\sum_{n\in\Z} g_{\gamma_n *}(u_{L_0,L_1,\gamma_n}(\mathrm{id}))
=\id\otimes \sum_{n\in\Z} q_e^{-n^2/2}\bigl(q_e^{n}t\bigr)^{n}
=\id\otimes \sum_{n\in\Z} q_e^{n^2/2}\, t^n
\]
（作为从 $\cO_X$ 到线丛 $\Phi(L_1)$ 的态射，在 $\Phi(L_1)$ 于 $X_K$ 的 $\Z$-覆盖 $K^*-q_e^\Z$ 上的平凡化中表出）；这就重新得到了椭圆曲线 $K^*/q_e^\Z$ 的经典 theta 函数。
\end{example}

现在说明 HPL 如何给出式~\eqref{eq:averaging} 中的平均。需要计算映射
\[
\Phi_{L,L'}\circ\tilde{i}^1:CF((L,\cE),(L',\cE'))\to \Hom(\Phi(L,\cE),\Phi(L',\cE')),
\]
其中 $\tilde{i}^1=(\id+h\delta^1)^{-1}i$，$i,h,\delta$ 如第~\ref{ss:HPLFloer} 节。取某 $x\in \cX(L,L';\varepsilon h)$ 与 $\varphi\in\Hom(\cE_x,\cE'_x)$，则 $\tilde{i}^1(\varphi)=\sum_{n=0}^{\infty}(-h\delta^1)^n i(\varphi)$。

映射 $h\delta^1:CF^0((L,\cE),(L',\cE');H)\to CF^0((L,\cE),(L',\cE');H)$ 可显式描述如下。公式~\eqref{eq:monomial_for_a_point} 给出等同
\[
CF^0((L,\cE),(L',\cE');H)\cong \bigoplus_{v\in V(G)}\Hom_{\cO_K}(\cE_v,\cE'_v)\oplus \bigoplus_{e\in E(G)}\Hom(\Phi(L,\cE)_{\mid U_e},\Phi(L',\cE')_{\mid U_e}).
\]
在此等同下，对 $\varphi\in \Hom_{\cO_K}(\cE_v,\cE'_v)$ 有
\[
h\delta^1(\varphi)=\sum_{e/v}(\varphi\otimes\cO_{U_v})_{\mid U_e},
\]
进一步，对每条边 $e:v\to v'$、$\varphi\in\Hom_{\cO_K}(\cE_v,\cE'_v)$ 与整数 $n$，传播规则给出：当 $n\geq r_e(L,L')$ 时
\[
h\delta^1\bigl((\varphi\cdot t_{e/v}^{-n})_{\mid U_e}\bigr)=\sum_{e'/v,\,e'\neq e}(\varphi\cdot t_{e/v}^{-n})_{\mid U_{e'}},
\]
而对 $0< n< r_e(L,L')$，
\[
h\delta^1\bigl((\varphi\cdot t_{e/v}^{-n})_{\mid U_e}\bigr)=\sum_{e'/v,\,e'\neq e}(\varphi\cdot t_{e/v}^{-n})_{\mid U_{e'}}
+\sum_{e''/v',\,e''\neq e}\bigl(R_{\cE',e/v}\varphi R_{\cE,e/v}^{-1}\bigr)\cdot\bigl(T^{-nS_e(L,L')}t_{e/v'}^{n-r_e(L,L')}\bigr)_{\mid U_{e''}}.
\]
由此即知 $\Phi_{L,L'}\circ\tilde{i}^1$ 恰给出式~\eqref{eq:averaging} 的平均。这里重要的是 $CF^{\bullet}(L,L')$ 集中在次数零，故所有旋转数 $r_e(L,L')$ 严格为正，传播时不会「丢掉」任何单项式。
% [Tutor Note: 易错] 上面的求和 $\sum_{n=0}^\infty(-h\delta^1)^n$ 在 $CF$ 完备化中才收敛；其每一项对应一个「经若干条边传播」的单项式，这正是 Schottky 群胚作用轨道的平均。丢了「$CF$ 集中在次数零」这一条，传播时就会出现负旋转数的项，平均式也不再收敛。依据：原文 §6 末段与本工程 03_术语表。
